% Unit 085 English translation of chapter6.tex, lines 2458--2603.
% Source: Methods of Algebra, Volume 2, commit
% 9a5803ff77dd3257484cb177f851a73770a59dd3 (CC BY 4.0).
% stable-unit-id: o014.aljabr2.chapter6.profinite-cohomology-b
% Coverage: the remainder of Section 6.12: delta functors, low-degree
% interpretations, derived-functor identification, profinite induction, and
% Shapiro's Lemma.

% segment-id: o014.li2.chapter6.profinite-cohomology-b.t001
\begin{theorem}\label{prop:profinite-cohomology-long}
	Let $G$ be a profinite group and let
	$0\to M'\xrightarrow{f}M\xrightarrow{f'}M''\to0$ be a short exact
	sequence of smooth $G$-modules. There is an associated long exact sequence
	\begin{multline*}
		\cdots\to\Hm^{n-1}(G,M'')\xrightarrow{\delta^{n-1}}
		\Hm^n(G,M')\xrightarrow{\Hm^n(f)}\Hm^n(G,M)\\
		\xrightarrow{\Hm^n(f')}\Hm^n(G,M'')\xrightarrow{\delta^n}
		\Hm^{n+1}(G,M')\to\cdots,
	\end{multline*}
% Editorial correction carried from the Indonesian witness: the Chinese source
% labels the middle arrow $\Hm^n(g)$ although the displayed map is $f'$.
	with the convention that $\Hm^n(G,\cdot)=0$ for $n<0$. The connecting
	homomorphisms $\delta^n$ are natural with respect to morphisms of short
	exact sequences; in other words, they form a cohomological $\delta$-functor.
\end{theorem}
\begin{proof}
	We work with the standard complex; the argument for the normalized version
	is similar. By Proposition~\ref{prop:long-exact-sequence-ses}, it suffices
	to prove that
	\[ 0\to C^n(G,M')\to C^n(G,M)\to C^n(G,M'')\to0 \]
	is a short exact sequence. The only point that is not immediate is the
	surjectivity of $C^n(G,M)\to C^n(G,M'')$. Let $u:G^n\to M''$ be
	continuous. As above, there is an open normal subgroup $K\lhd G$ such
	that $u$ is constant on each coset of $K^n\lhd G^n$. On every coset,
	choose an arbitrary inverse image in $M$ of this constant value.
\end{proof}

% segment-id: o014.li2.chapter6.profinite-cohomology-b.p001
Let $M$ be a smooth $G$-module with the discrete topology. The low-degree
terms of the complex $C(G,M)$ or $\overline{C}(G,M)$ still have direct
interpretations.
\begin{description}
	\item[($n=0$)] Clearly $C^0(G,M)=M$ and
	$\Hm^0(G,M)=Z^0(G,M)=M^G$. The normalized standard complex gives the
	same result.
	\item[($n=1$)] The elements of
	$Z^1(G,M)=\overline{Z}^1(G,M)$ are continuous crossed homomorphisms
	$G\to M$. Equivalently, they are continuous homomorphisms
	$\sigma:G\to M\rtimes G$ satisfying $\pi\sigma=\identity_G$, as in
	Proposition~\ref{prop:crossed-semidirect-product}. Here
	$\pi:M\rtimes G\to G$ is the projection and $M\rtimes G$ carries the
	product topology. On the other hand, the elements of $B^1(G,M)$ are
	generated by the crossed homomorphisms $f_m:g\mapsto gm-m$; these are
	automatically continuous.

	\item[($n=2$)] Suppose that $M$ is finite. The elements of
	\[ \Hm^2(G,M)\simeq\frac{Z^2(G,M)}{B^2(G,M)}
	\simeq\frac{\overline{Z}^2(G,M)}{\overline{B}^2(G,M)} \]
	correspond bijectively to equivalence classes of extensions of profinite
	groups
	\[ 1\to M\to E\xrightarrow{\pi}G\to1
	\qquad\text{(short exact sequence of groups)}, \]
% Editorial correction carried from the Indonesian witness: group extensions
% have terminal identity groups $1$, not zero groups as printed in the source.
	where the conjugation action of $G$ on $M$ comes from the $G$-module
	structure on $M$. Compared with the version in
	Definition~\ref{def:group-ext} and \eqref{eqn:group-ext-diagram}, a
	``profinite group extension'' and its classification must additionally
	satisfy:
	\begin{compactitem}
		\item $E$ is a profinite group, $\pi$ is continuous, and $M$ is
		embedded as a closed subgroup of $E$;
		\item equivalences, or isomorphisms, of extensions are defined by
		continuous group homomorphisms $E\to E'$, and the splitting of an
		extension is likewise defined topologically.
	\end{compactitem}

	Moreover, the automorphism group of an extension is isomorphic to
	$\overline{Z}^1(G,M)$. The proofs are not difficult, but require some
	topological language and are therefore left as exercises for this chapter;
	the finiteness of $M$ is essential.
\end{description}

% segment-id: o014.li2.chapter6.profinite-cohomology-b.p002
We now return to the right derived functors
$(\mathrm{R}^n(\cdot)^G)_{n\geq0}$. One could study their values on arbitrary
bounded-below complexes\footnote{In other words, study the hypercohomology of
profinite groups.}, or even their operation in the derived category. This
section, however, discusses only smooth $G$-modules; the rest is left to the
exercises.

% segment-id: o014.li2.chapter6.profinite-cohomology-b.t002
\begin{proposition}\label{prop:profinite-cohomology-identification}
	There is an isomorphism of cohomological $\delta$-functors
	\[ (\Hm^n(G,\cdot))_{n\geq0}
	\simeq(\mathrm{R}^n(\cdot)^G)_{n\geq0}. \]
\end{proposition}
\begin{proof}
	Theorem~\ref{prop:profinite-cohomology-long} shows that both sides are
	cohomological $\delta$-functors, and at $n=0$ both are the invariants
	functor $(\cdot)^G$. By Proposition~\ref{prop:erasable-univ-delta}, it is
	enough to show that $\Hm^n(G,\cdot)$ is erasable for $n>0$ in the sense
	of Definition~\ref{def:erasable}. Let $M$ be a smooth $G$-module and embed
	it in an injective smooth $G$-module $I$. For every open normal subgroup
	$K\lhd G$, the adjunction \eqref{eqn:invariant-adjunction-profinite}
	shows that $(\cdot)^K$ preserves injective objects. Hence $I^K$ is an
	injective $G/K$-module, and
	\[ n>0\implies\Hm^n(G,I)
	=\varinjlim_K\Hm^n(G/K,I^K)=0. \]
	This proves erasability.
\end{proof}

% segment-id: o014.li2.chapter6.profinite-cohomology-b.p003
By Proposition~\ref{prop:profinite-cohomology-identification}, the method of
\S\ref{sec:change-of-group} applies directly to a continuous homomorphism
$\varphi:H\to G$ of profinite groups and defines a cohomological
``pullback'' operation
\[ \Hm^n(G,M)\to\Hm^n(H,\varphi^*M),\qquad
	M:\text{ a smooth $G$-module},\quad n\in\Z_{\geq0}. \]
These operations are compatible with long exact sequences. As usual, exactness
of $\varphi^*$ and the universal property of cohomological $\delta$-functors
reduce the matter to the evident inclusion at $n=0$,
$M^G\hookrightarrow(\varphi^*M)^H$.

% segment-id: o014.li2.chapter6.profinite-cohomology-b.p004
As a special case, for all $n$ and all closed subgroups $H$ there are
restriction maps
$\mathrm{res}^n:\Hm^n(G,M)\to\Hm^n(H,M)$; here
$\Res^G_HM$ is abbreviated to $M$.

% segment-id: o014.li2.chapter6.profinite-cohomology-b.p005
The homomorphisms above may also be constructed from the homomorphisms
\begin{align*}
	\Hm^n(G/K,M^K)&\xrightarrow{\eqref{eqn:change-of-group-coh}}
	\Hm^n(H/\varphi^{-1}(K),\varphi_K^*(M^K))\\
	&\xrightarrow{M^K\subset(\varphi^*M)^{\varphi^{-1}K}}
	\Hm^n(H/\varphi^{-1}(K),(\varphi^*M)^{\varphi^{-1}K}),
\end{align*}
where $K$ ranges over the open normal subgroups of $G$. Reducing to the case
of finite groups shows that the homomorphism on the standard complex is
computed by the formula of Proposition~\ref{prop:grp-equiv}. To prove that
this map is indeed the same as the one constructed earlier by the universal
property, it is enough to compare them at $n=0$ and show compatibility with
long exact sequences. The standard complex suffices for both tasks.

% segment-id: o014.li2.chapter6.profinite-cohomology-b.p006
Similarly, the corestriction homomorphism of
Definition--Proposition~\ref{def:cor} extends to profinite groups. Let $H$
be an open subgroup of $G$. It is then automatically closed and $(G:H)$ is
finite (see \S\ref{sec:profinite-groups}). Define homomorphisms
\[ \mathrm{cor}^n:\Hm^n(H,M)\to\Hm^n(G,M),\qquad
	M:\text{ a smooth $G$-module},\quad n\in\Z_{\geq0}, \]
which are compatible with long exact sequences and at $n=0$ give
$\nu^{G|H}:M^H\to M^G$.

% segment-id: o014.li2.chapter6.profinite-cohomology-b.p007
The construction again uses the universal property of cohomological
$\delta$-functors, but here we need to know that
$\Res^G_H:G\dcate{Mod}^\infty\to H\dcate{Mod}^\infty$ preserves injective
objects; this is the content of
Corollary~\ref{prop:Res-injective-profinite}. Once this fact is granted,
Proposition~\ref{prop:res-cor} still has an analogue: for every $n$,
\[ \mathrm{cor}^n\circ\mathrm{res}^n
	=(G:H)\cdot\identity_{\Hm^n(G,M)}. \]
As before, the proof reduces to the case $n=0$.

% segment-id: o014.li2.chapter6.profinite-cohomology-b.p008
We now discuss induced modules in the profinite case. Henceforth $G$ is
always a profinite group. To distinguish the two settings, denote the
induction functor without topology (Definition~\ref{def:induced-module}) by
$\Ind^{G,\mathrm{alg}}_H$; following Lemma~\ref{prop:Ind-as-mappings}, we
realize it as a mapping space.

% segment-id: o014.li2.chapter6.profinite-cohomology-b.d001
\begin{definition}
	\index{$G$-module!induced}
	For a closed subgroup $H$ of $G$ and a smooth $H$-module $N$, equip the
	mapping space $\mathrm{Maps}(G,N):=\{f:G\to N\}$ with pointwise addition
	and scalar multiplication, together with the left $G$-action
	\[ (gf)(x)=f(xg),\qquad g,x\in G. \]
	Define the smooth $G$-module
	\begin{align*}
		\Ind^G_H(N)&:=\{f\in\mathrm{Maps}(G,N)^\infty:
		\forall h\in H,\ \forall x\in G,\ f(hx)=hf(x)\}\\
		&=(\Ind^{G,\mathrm{alg}}_H(N))^\infty.
	\end{align*}
	The notation $(\cdot)^\infty$ is from
	Definition--Proposition~\ref{def:smooth-part}. This is called the
	$G$-module induced from $N$.
\end{definition}

% segment-id: o014.li2.chapter6.profinite-cohomology-b.p009
Since $G$ is compact, smoothness of
$f\in\Ind^{G,\mathrm{alg}}_H(N)$ is equivalent to its continuity as a map
$G\to N$, where $N$ carries the discrete topology.

% segment-id: o014.li2.chapter6.profinite-cohomology-b.t003
\begin{proposition}\label{prop:Ind-injective-profinite}
	For every closed subgroup $H$ of $G$, there is an adjoint pair
	\[ \Res^G_H:G\dcate{Mod}^\infty
	\leftrightarrows H\dcate{Mod}^\infty:\Ind^G_H. \]
	Consequently, $\Ind^G_H$ preserves injective objects.

	If $H$ is an open subgroup of $G$, there is also an adjoint pair
	\[ \Ind^G_H:H\dcate{Mod}^\infty
	\leftrightarrows G\dcate{Mod}^\infty:\Res^G_H. \]
	In this case $\Ind^G_H$ also preserves projective objects.
\end{proposition}
\begin{proof}
	By the adjunction in the setting without topology and
	Definition--Proposition~\ref{def:smooth-part}, for every smooth
	$G$-module $M$ and smooth $H$-module $N$ there are canonical
	isomorphisms
	\begin{multline*}
		\Hom_H(\Res^G_H(M),N)
		\simeq\Hom_G(M,\Ind^{G,\mathrm{alg}}_H(N))\\
		=\Hom_G(M,\Ind^{G,\mathrm{alg}}_H(N)^\infty)
		=\Hom_G(M,\Ind^G_H(N)).
	\end{multline*}

	Now suppose that $H$ is open, so that $(G:H)$ is finite. From
	Corollary~\ref{prop:iInd-Ind} we obtain the adjunction without topology
	\[ \Hom_H(N,\Res^G_H(M))
	\simeq\Hom_G(\Ind^{G,\mathrm{alg}}_H(N),M). \]
	We claim that
	$\Ind^{G,\mathrm{alg}}_H(N)=\Ind^G_H(N)$. Take
	$f\in\Ind^{G,\mathrm{alg}}_H(N)$ and $x\in G$, and choose an open
	subgroup $K_0\subset H$ such that $f(x)\in N^{K_0}$. Set
	$K:=x^{-1}K_0x$, which is an open subgroup of $G$. Then
	\[ g\in K\implies xg=xgx^{-1}x\in K_0x
	\implies f(xg)=f(x). \]
	Since $x$ was arbitrary, the map $f:G\to N$ is locally constant, so
	$f\in\Ind^G_H(N)$. This proves the second adjunction.

	Finally, the claims about preserving injective or projective objects
	follow from exactness of $\Res^G_H$ and
	Proposition~\ref{prop:adjoint-injective-projective}.
\end{proof}

% segment-id: o014.li2.chapter6.profinite-cohomology-b.l001
\begin{lemma}\label{prop:Ind-exact-profinite}
	If $H$ is a closed subgroup of $G$, then
	$\Ind^G_H:H\dcate{Mod}^\infty\to G\dcate{Mod}^\infty$ is exact.
\end{lemma}
\begin{proof}
	We already know that $\Ind^G_H$ has a left adjoint, so it is left exact.
	It remains to prove that $N_1\twoheadrightarrow N_2$ implies
	$\Ind^G_H(N_1)\twoheadrightarrow\Ind^G_H(N_2)$. By
	Lemma~\ref{prop:profinite-section}, the quotient homomorphism
	$\pi:G\to H\backslash G$ has a continuous section $s$. For every smooth
	$H$-module $N$, there is a $\Bbbk$-module isomorphism
	\[\begin{tikzcd}[row sep=tiny]
		\Ind^G_H(N) \arrow[leftrightarrow,r,"\sim"] &
		\{\text{continuous maps }H\backslash G\to N\}\\
		f \arrow[mapsto,r] & {[x\mapsto f(s(x))]}\\
		{[hs(x)\mapsto hf'(x)]} & f' \arrow[mapsto,l]
	\end{tikzcd}\]
	where $x\in H\backslash G$ and $h\in H$. Although this isomorphism does
	not preserve the $G$-action, it is natural in $N$.

	Since continuous maps $H\backslash G\to N$ are always locally constant,
	working on the right-hand side immediately shows that
	$\Ind^G_H(N_1)\to\Ind^G_H(N_2)$ is surjective.
\end{proof}

% segment-id: o014.li2.chapter6.profinite-cohomology-b.t004
\begin{corollary}\label{prop:Res-injective-profinite}
	If $H$ is an open subgroup of $G$, then
	$\Res^G_H:G\dcate{Mod}^\infty\to H\dcate{Mod}^\infty$ preserves
	injective objects.
\end{corollary}
\begin{proof}
	By the preceding two results, it has the exact left adjoint $\Ind^G_H$.
\end{proof}

% segment-id: o014.li2.chapter6.profinite-cohomology-b.p010
We now obtain the profinite version of Theorem~\ref{prop:Shapiro}.

% segment-id: o014.li2.chapter6.profinite-cohomology-b.t005
\begin{proposition}
	\index{Shapiro's lemma}
	Let $H$ be a closed subgroup of a profinite group $G$ and let $N$ be a
	smooth $H$-module. There are isomorphisms
	\[ \Hm^n(G,\Ind^G_H(N))\simeq\Hm^n(H,N),
	\qquad n\in\Z_{\geq0}. \]
\end{proposition}
\begin{proof}
	Because $\Ind^G_H$ is exact and preserves injective objects
	(Lemma~\ref{prop:Ind-exact-profinite} and
	Proposition~\ref{prop:Ind-injective-profinite}), the argument in
	Theorem~\ref{prop:Shapiro} using erasable cohomological $\delta$-functors
	applies unchanged. Recall that the heart of that argument is the case
	$n=0$, namely
	\[ \Ind^G_H(N)^G
	=(\Ind^{G,\mathrm{alg}}_H(N)^\infty)^G
	=(\Ind^{G,\mathrm{alg}}_H(N))^G\rightiso N^H. \]
	The first equality is the definition, the second is immediate, and the
	last isomorphism comes from Corollary~\ref{prop:Ind-invariant}.
\end{proof}

% segment-id: o014.li2.chapter6.profinite-cohomology-b.p011
In addition:
\begin{itemize}
	\item the cohomological Lyndon--Hochschild--Serre spectral sequence
	(see \S\ref{sec:LHS-SS}) extends to the case where $G$ is a profinite
	group and $H$ is a closed normal subgroup;
	\item the cup product (see \S\ref{sec:cup-product}) also extends to
	profinite groups by using the direct definition on the standard complex;
	see Definition--Proposition~\ref{def:group-cup-3}.
	\index{cup product}
\end{itemize}
The statements are unchanged, so we do not repeat them.
